\ficha{Marcondes (1998), Padrão-ouro e estabilidade}{marcondes1998}

\begin{identificacao}
MARCONDES, Renato Leite. Padrão-ouro e estabilidade. \textit{Estudos Econômicos},
São Paulo, v. 28, n. 3, p. 533-559, jul./set. 1998.

Artigo de revisão da literatura, 27 páginas, lido na íntegra, em português.
Versão lida idêntica à publicada.
\end{identificacao}

\bloco{Problema e tese}
Por que o mesmo arranjo monetário produziu estabilidade entre 1879 e 1913 e
instabilidade entre 1919 e 1939? A pergunta serve para avaliar a proposta,
corrente nos anos 1980 e 1990, de restaurar um lastro metálico como âncora
nominal. A resposta tem duas camadas. A diferença entre os períodos não está no
automatismo do sistema, mas em duas propriedades externas a ele, cooperação entre
bancos centrais e credibilidade do compromisso de cada país com a paridade,
presentes antes de 1914 e ausentes depois. E a própria estabilidade do período
áureo é em boa parte mito estatístico: o ouro entregou preços estáveis no longo
prazo ao custo de renda instável e de forte transmissão de choques.

\bloco{Argumento}
\begin{enumerate}
\item Contra Kindleberger (1973), que atribui o sucesso pré-guerra à Inglaterra
como emprestadora de última instância, Eichengreen mostra a atuação do Banco da
França e do Reichsbank nas crises de 1890 e 1907 (p. 535). No entreguerras não
faltou hegemonia, faltaram credibilidade e cooperação (p. 536).

\item As regras do jogo são conversibilidade a paridade fixa, liberdade de
importar e exportar ouro e vínculo entre meio circulante e estoque metálico
(p. 540). Delas decorreriam o ajuste automático de Hume e a simetria de preços.
Na prática o ajuste passou pela taxa de juros e pela esterilização dos fluxos, e
McCloskey e Zecher (1976) acham relação empírica fraca entre fluxos de ouro e
crédito interno (p. 541, nota 4).

\item A assimetria é o ponto mais útil ao capítulo. Na expansão, preços altos de
commodities e crédito abundante acomodam a periferia; na contração, um pequeno
aumento da taxa de desconto do Banco da Inglaterra repatria capitais, derruba os
preços das matérias-primas e agrava o serviço da dívida (p. 542).

\chave{Havia uma assimetria de ajustamento do sistema, os países centrais
apresentavam condições facilitadas, enquanto os países periféricos sofriam sérias
dificuldades em permanecer no padrão-ouro nas fases contracionistas.}{p. 542}

\item No entreguerras a restauração é desencontrada, paridade pré-guerra na
Inglaterra em 1925 e paridade desvalorizada na França em 1928, com reservas
heterogêneas sob o gold exchange standard saído de Gênova (p. 545-546). Somam-se
o deslocamento da política econômica para renda e emprego (p. 547-548) e a
fragmentação federativa do Federal Reserve, que reduziu a disposição americana de
cooperar (p. 548). Na depressão o padrão-ouro opera como propagador, e o abandono
precoce pela Inglaterra em 1931 explica sua recuperação anterior à americana
(p. 550).

\item A seção comparativa desmonta a premissa do debate. Schwartz (1987), Meltzer
e Robinson (1989) e Bordo (1993) convergem: inflação média mais baixa e
persistência inflacionária menor sob o ouro, mas variabilidade maior de preços e
de produto, e melhor desempenho do Bretton Woods conversível (p. 552-555). Bordo
credita a durabilidade do ouro à flexibilidade de preços e salários e ao pequeno
compromisso com o pleno emprego, e julga isso insuficiente sem credibilidade
(p. 555).
\end{enumerate}

\chave{a evidência empírica, como vimos acima em Schwartz, Meltzer e Robinson e
Bordo, não corrobora tal apregoada estabilidade}{p. 556}

\bloco{Conceitos e definições operacionais}
Padrão-ouro é definido pelas três regras do jogo da p. 540, não pela existência de
lastro. Credibilidade é a convicção do público de que a obrigação de manter a
paridade continuará, e é ela que gera a automaticidade do ajuste, dispensando
decisão nova a cada flutuação (p. 543). Cooperação é o auxílio mútuo entre bancos
centrais, por acordo tácito ou coordenado (p. 535). O sistema é lido como regra
contra discricionariedade, \textit{rules versus discretion} (p. 536). Âncora
nominal é a determinação endógena do nível de preços pela oferta e pela demanda de
ouro (p. 540). A cláusula de escape aparece sem esse nome: suspender a
conversibilidade por guerra ou choque grande é admissível desde que o país se
comprometa a voltar à paridade anterior, ainda que ao custo de recessão
deflacionária (p. 540).

\bloco{Evidência e método}
Resenha crítica, sem fonte primária e sem estimação própria. O material empírico é
reproduzido de terceiros: crescimento de M1 por região entre 1926 e 1929, de
Eichengreen (p. 549); inflação, renda real per capita, desemprego e oferta de
moeda no Reino Unido e nos Estados Unidos, de Schwartz, com 1870 a 1914 contra
1946 a 1979 (p. 553); e estatísticas descritivas de sete países entre 1881 e 1989
por regime, de Bordo, onde a renda per capita cresce 1,5 com desvio padrão 3,7 sob
o ouro contra 4,2 e 2,7 sob Bretton Woods (p. 554). Não há contrafactual formal.

\bloco{Agentes e vertentes}
Eichengreen (1985, 1992) é a espinha dorsal, grafado ``Einchengreen'' no corpo, e
fornece a crítica à estabilidade hegemônica e o par cooperação e credibilidade.
Kindleberger (1973) ocupa a posição contrária, da liderança hegemônica, e é
derrotado na exposição. Triffin (1985) sustenta a assimetria centro-periferia, com
Nurkse (1954) nos fluxos de capital. A vertente monetarista entra em bloco:
Friedman (1992) no bimetalismo, Friedman e Schwartz (1963), Bordo (1989) e Lucas
(1994) na depressão como falha do Federal Reserve, com Temin (1976) registrado
como dissenso (p. 549, nota 7). A avaliação comparada de regimes reúne Schwartz
(1987), Meltzer e Robinson (1989) e Bordo (1993), os três desfavoráveis à
superioridade do ouro. Somam-se McCloskey e Zecher (1976) na abordagem monetária
do balanço de pagamentos, Huffman e Lothian (1984) na transmissão de ciclos e
Crabbe (1989) no Federal Reserve. Keynes é a voz crítica de época, com a relíquia
bárbara de 1923 (p. 546), e Polanyi, grafado ``Polany'', fecha o entreguerras com
a leitura institucional do fim da sociedade de mercado do século XIX (p. 552).

\bloco{Críticas e limites}
Não é pesquisa original e não produz evidência nova. Descreve a cláusula de escape
mas ignora a literatura que a formalizou, Bordo e Kydland (1990) e Bordo e
Schwartz (1994), e desconhece o argumento de reputação e acesso a crédito de Bordo
e Rockoff (1996), justamente a parte da bibliografia que mais interessa a um país
periférico. Fica sem solução a tensão entre creditar o sucesso pré-guerra à
credibilidade e aceitar que esse período foi o mais instável em renda. A periferia
aparece só como receptora de choque, sem discussão de por que esses países aderiam
nem do prêmio de risco que pagavam. O Brasil não é objeto: é citado uma vez, como
caso de inflação crônica nos anos 1980 (p. 534), e a América Latina só como
tomadora de empréstimos no colapso de 1929 a 1931 (p. 548-549).

\begin{dialogo}
\item Regra contra arbítrio. Procurar em O Paiz, no Correio da Manhã e nos
Documentos Parlamentares de 1906 a defesa da Caixa como mecanismo que retira do
governo a decisão sobre a taxa, com vocabulário de automatismo e de fim do arbítrio
do Tesouro. Testa se a oposição \textit{rules versus discretion} (p. 536) é
categoria nativa do debate ou importação do analista.
\item Assimetria do ciclo do centro. Procurar, em 1907 e 1908 e sobretudo em 1913
e 1914, menções à taxa de desconto do Banco da Inglaterra, à retirada de capitais
e à queda do preço do café como causas da pressão sobre a Caixa. Testa a hipótese
de Triffin retomada na p. 542.
\item Substituto da cooperação. O Brasil não participa da rede de bancos centrais
a que Marcondes credita a estabilidade. Procurar como os jornais tratam a agência
de Londres e a custódia do ouro em 1911, se como garantia técnica, se como tutela
dos banqueiros, e verificar se o arranjo com os Rothschild cumpre ali essa função.
\item Suspensão como cláusula de escape. Procurar, em agosto e setembro de 1914,
os termos que justificam a suspensão do troco, se abandono do sistema, se
interrupção temporária com promessa de retorno à mesma paridade (p. 540).
\item Objetivo interno contra externo. Procurar, no debate sobre 12, 15 ou 27
pence, o argumento de que sustentar a paridade sacrifica renda, emprego e a
lavoura. Marcondes situa esse tipo de argumento no entreguerras (p. 547-548), de
modo que achá-lo em 1906 antecipa o deslocamento e é achado, não confirmação.
\end{dialogo}

\usonocapitulo{Parte I. Mapa em português da bibliografia anglófona sobre por que
o padrão-ouro clássico funcionou. Sustenta duas afirmações do capítulo: que a
estabilidade pré-guerra dependeu de credibilidade e cooperação, não do automatismo
do mecanismo de Hume, e que a assimetria de ajustamento entre centro e periferia é
estrutural. Serve ainda para relativizar a estabilidade atribuída ao sistema. Não
trata do Brasil nem de caixas de conversão, e portanto não sustenta afirmação
alguma sobre a Caixa: entra como enquadramento e referência secundária,
subordinada a Eichengreen, Bordo e Triffin no original. A frase que fecha o artigo
desloca o problema do lastro para as instituições que sustentam o compromisso.}

\chave{Se tais propriedades estão presentes, possivelmente não haveria necessidade
do padrão-ouro para conseguirmos estabilidade}{p. 557}
